\scanpage{39}
% archive foliation: 38

\uline{Récapitulation partielle}

A un module $M$ libre de rang $2$ sur $\mathbb{Z}/4\mathbb{Z}$ on associe

\[
(49)\quad
\left\{
\diagram{\begin{array}{p{.88\linewidth}}
a) L'octaèdre (non orienté) \inserted{de sommets} $\Delta_M=M^*/\pm1$, l'ens des trois diagonales s'identifiant à $(M/2M)^*\simeq M_0^*=J$\\[1em]
b) L'ensemble à quatre éléments $I$, qu'on peut p.ex. déduire de $\Delta$ comme
\[
I_M\simeq
\underset{\substack{\text{faces de }\Delta\\\unsure{\text{ou sections de }\Delta\text{ sur }J}}}
{\underbrace{F(\Delta)}}/\text{antipodisme}
\]\\
c) Le carré $Q_M$ \inserted{d'ens de sommets} $\unsure{S_M}=P_M/M_0$, où $P_M$ est le $\mathrm{End}(M_0)$-torseur formé des sections de $M^*$ sur $M_0^*$ satisfaisant $\xi_i+\xi_j+\xi_k=0$, et où on identifie $\mathrm{End}(M_0):M_0\oplus\mathbb{F}_2^{\underline\omega}$, de sorte que $P_M/M_0$ est un $\mathbb{F}_2^{\underline\omega}$ torseur, donc est l'ens des sommets d'un carré \inserted{$Q=Q_M$} dont l'ens des codiagonales est $\underline\omega$.
\end{array}}
\right.
\]

\margin{On a un iso. can.
$J\simeq\mathcal P_{2,2}(I_M)$,
$I_M$ torseur sous $\underset{(\simeq\mathbb{F}_2^J/\mathbb{F}_2)}{\underbrace{V_J(\mathbb{F}_2)\simeq M_0}}$
$I_M\simeq P_M/\mathbb{F}_2^{\underline\omega}$.}

\margin{\uline{NB} $F(\Delta)\simeq I\times
\underset{\text{``couleurs'' des faces}}{\underset{\downarrow}{C_\Delta}}$}

\margin{où $\underline\omega\simeq\underline\omega_I\simeq\underline\omega_J$}

Les données a) et b) ne dépendent que de la \deleted{fibré} droite projective $X_M=\mathbb{P}(M)$ sur $\mathbb{Z}/4\mathbb{Z}$ ($\Delta$ est l'ens des p$^{\mathrm{ts}}$ de celle-ci, $J$ celui des p$^{\mathrm{ts}}$ de $(X_M)_0$, droite projective sur $\mathbb{F}_2$), mais non pas $Q_M$. On a les bijections canoniques
\[
(50)\quad
\left\{
\begin{array}{lll}
\underset{\text{codiag de }Q}{\underbrace{C_Q}}
\simeq
\underset{\deleted{\unsure{\text{orientations de }\Delta}}}{\underline\omega_\Delta}
\overset{\mathrm{déf}}{\simeq}\underline\omega_I\simeq\underline\omega_J
&\text{soit }\underline\omega\\
\underset{\text{diag de }Q}{D_Q}
\simeq
\underset{\text{``couleurs'' des faces}}{\underline C_\Delta}
\simeq\underline\omega\wedge\underline\mu
&\text{où }\underline\mu\text{ défini ci-dessous}\\
\underset{\text{or. de }Q}{\underline\omega_Q}
\simeq\underline{\mathrm{Or}}(\Delta)\simeq\det(M)^*
&\text{soit }\underline\mu
\end{array}
\right.
\]
